\section{Geometría orbital 3D}

\subsection{Visualización del sistema perifocal}

Mostrar en 3D los vectores $\hat{p}$ (hacia el perihelio), $\hat{q}$ (perpendicular dentro del plano orbital) y $\hat{w} = \vec{h}/h$ (normal al plano) para la órbita de Apophis. Relacionar estos vectores con los ángulos $\Omega$, $I$ y $\omega$, y con el sistema de referencia eclíptico.

\subsection{Plano orbital vs. plano de la eclíptica}

Construir una visualización 3D que muestre el plano orbital de Apophis inclinado $3.34^{\circ}$ respecto al plano de la eclíptica, junto con la línea de nodos como intersección entre ambos planos. Mostrar a la Tierra cruzando esa línea durante el encuentro de 2029.

\subsection{Aplicación de las matrices de Euler}

Tomar las coordenadas en el sistema perifocal, $\vec{r}_{pf} = (r \cos f, r \sin f, 0)$, y aplicar explícitamente la transformación
\[
\vec{r}_{obs} = R_z(\Omega)\, R_x(I)\, R_z(\omega)\, \vec{r}_{pf}.
\]
Verificar que el resultado coincide con las coordenadas entregadas por \texttt{REBOUND}.